\section{Honest intervals under rare outcome arrival}

The deterministic squared-error envelope in \cref{obj:thm:unrestricted-mixed-count-upper} also supports uncertainty quantification for the average treatment effect. The relevant criterion combines uniform coverage over the rare-arrival model with maximal expected interval length. Fix a miscoverage level \leanref{sym:\alpha}{\ensuremath{\alpha\in(0,1)}}, so that the required coverage is \(1-\alpha\). To include procedures with auxiliary randomization, both coverage and expected length are evaluated under the joint distribution of the observed sample and the reported endpoints. The following definition specifies this decision problem.

\begin{definitionv}{def:interval-length-risk}[Minimax expected interval length]
The minimax maximal expected interval length, \(\mathfrak L_{n,d,q,\alpha}\), is defined for \(n,d\in\mathbb N\) and \(q,\alpha\in\mathbb R\) as follows. Let \(\mathcal M_{n,d,q}\) be the model class in \cref{obj:def:model-class}, let \(\mathbf O_n\) be the observed sample with law \(\mathcal L_P(\mathbf O_n)\) from \cref{obj:def:observed-experiment}, and let \(\tau(P)\) be the average treatment effect in \cref{obj:def:ate-functional}.

The class \(\mathcal C_{n,d,q,\alpha}\) consists of parameter-independent Markov kernels \(\mathsf I\) from the observed-sample space to endpoint pairs \((l,u)\) satisfying \(-1\le l\le u\le1\), with uniform joint sample-and-randomization coverage: for every \(P\in\mathcal M_{n,d,q}\),
\[
\int\!\int
\mathbf1\{l\le\tau(P)\le u\}\,
\mathsf I(o,d(l,u))\,
\mathcal L_P(\mathbf O_n)(do)
\ge 1-\alpha.
\]
The realized connected interval is \([l,u]\), with length \(u-l\). Define
\[
\mathfrak L_{n,d,q,\alpha}
:=
\inf_{\mathsf I\in\mathcal C_{n,d,q,\alpha}}
\sup_{P\in\mathcal M_{n,d,q}}
\int\!\int (u-l)\,
\mathsf I(o,d(l,u))\,
\mathcal L_P(\mathbf O_n)(do),
\]
where the infimum and supremum are real-valued, with the convention that an infimum or supremum over an empty set equals zero. In particular, \(\mathfrak L_{n,d,q,\alpha}=0\) when \(\mathcal C_{n,d,q,\alpha}\) is empty. Both coverage and expected length integrate over the sample and the kernel draw. Deterministic connected-interval procedures are included through Dirac kernels.
\end{definitionv}

Uniform honesty requires the coverage inequality separately for every law in \cref{obj:def:model-class}. For a randomized endpoint procedure \leanref{sym:\mathsf I}{\ensuremath{\mathsf I}}, this requirement averages over its auxiliary draw as well as sampling variation. The same averaging governs length. Thus the admissible class \leanref{sym:\mathcal C_{n,d,q,\alpha}}{\ensuremath{\mathcal C_{n,d,q,\alpha}}} treats deterministic and randomized procedures within one criterion, and the minimax length \leanref{sym:\mathfrak L_{n,d,q,\alpha}}{\ensuremath{\mathfrak L_{n,d,q,\alpha}}} measures the smallest attainable worst-case expected width at the specified coverage.

A concrete honest interval uses the mixed-count estimator \leanref{sym:\widehat\tau^{\mathrm{mix}}_{n,d,q}}{\ensuremath{\widehat\tau^{\mathrm{mix}}_{n,d,q}}} from \cref{obj:def:mixed-count-estimator} and its deterministic risk envelope \leanref{sym:\mathfrak U_{n,d,q}}{\ensuremath{\mathfrak U_{n,d,q}}} in \cref{obj:thm:unrestricted-mixed-count-upper}. Both are calibrated to the supplied arrival floor \(q\). The envelope determines a radius before observing the sample; intersecting the resulting interval with the target range gives the following construction. Its coverage guarantee therefore requires the supplied floor to be no larger than every occupied-cell arrival probability.

\begin{definitionv}{def:risk-envelope-interval}[Risk-envelope interval \(I^{\mathrm{mix}}_\alpha\)]
The risk-envelope interval is
\[
I^{\mathrm{mix}}_\alpha
:=
\left[
\widehat\tau^{\mathrm{mix}}_{n,d,q}
-
\sqrt{\frac{\mathfrak U_{n,d,q}}{\alpha}},
\,
\widehat\tau^{\mathrm{mix}}_{n,d,q}
+
\sqrt{\frac{\mathfrak U_{n,d,q}}{\alpha}}
\right]\cap[-1,1],
\]
where \(\widehat\tau^{\mathrm{mix}}_{n,d,q}\) is the mixed-count estimator and \(\mathfrak U_{n,d,q}\) is its explicit deterministic squared-error envelope.
\end{definitionv}

The interval \leanref{sym:I^{\mathrm{mix}}_\alpha}{\ensuremath{I^{\mathrm{mix}}_\alpha}} is deterministic conditional on the observed sample and therefore enters \cref{obj:def:interval-length-risk} through a Dirac endpoint rule. Its radius converts uniform squared-error control into coverage at least \(1-\alpha\). The next result compares its maximal expected length with the best uniformly honest connected interval, using the comparison rate \leanref{sym:r_{n,d,q}}{\ensuremath{r_{n,d,q}}} defined in \cref{obj:def:frontier-rate}.

\begin{theoremv}{thm:connected-interval-frontier}[Connected interval length frontier]
For every miscoverage level \(\alpha\in(0,1)\), there exist constants \(0<\underline c_\alpha\le\overline C_\alpha<\infty\) such that the following holds for every \(n,d\in\mathbb N\) and \(q\in\mathbb R\) satisfying:
\begin{itemize}
\item \textbf{(Sample size and dimension.)} \(n\ge1\) and \(d\ge1\).
\item \textbf{(Arrival floor.)} \(0<q\le1\).
\item \textbf{(Rare-arrival slice.)} The restriction in \cref{obj:ass:rare-arrival-slice} holds.
\end{itemize}
Let \(r_{n,d,q}\) be the frontier rate in \cref{obj:def:frontier-rate}, let \(\mathfrak L_{n,d,q,\alpha}\) be the minimax expected interval length in \cref{obj:def:interval-length-risk}, and let \(I^{\mathrm{mix}}_\alpha\) be the interval in \cref{obj:def:risk-envelope-interval}. Over the model class \(\mathcal M_{n,d,q}\) in \cref{obj:def:model-class},
\[
\underline c_\alpha\sqrt{r_{n,d,q}}
\le \mathfrak L_{n,d,q,\alpha}
\le \sup_{P\in\mathcal M_{n,d,q}}
E_P\!\left[\operatorname{length}(I^{\mathrm{mix}}_\alpha)\right]
\le \overline C_\alpha\sqrt{r_{n,d,q}}.
\]
Here expectations and probabilities are taken under the law of the independent observed sample induced by \(P\), and interval length is defined by
\[
\operatorname{length}([l,u])=u-l.
\]
Moreover, for every \(P\in\mathcal M_{n,d,q}\), the average treatment effect \(\tau(P)\) from \cref{obj:def:ate-functional} satisfies
\[
\Pr_P\!\left\{\tau(P)\in I^{\mathrm{mix}}_\alpha\right\}\ge1-\alpha.
\]
\end{theoremv}

The upper bound reflects the square-root conversion from squared-error risk to an interval radius: the deterministic envelope controls coverage, while clipping bounds the realized length by the width of the untruncated interval. The converse connects statistical indistinguishability with the geometry of a connected interval. Coverage transfers between competing laws in the joint sample-and-endpoint distribution, and a realized interval that meets several separated target neighborhoods must span their separation. Working with this joint distribution retains the auxiliary endpoint draw throughout both the coverage comparison and the length calculation. The parameter-independent decision-rule comparison in \cref{obj:lem:randomized-kernel-tv-comparison} supplies the corresponding control when an endpoint procedure is appended to the observed experiment.

For each fixed coverage level \(1-\alpha\in(0,1)\), \cref{obj:thm:connected-interval-frontier} establishes expected-length order \(\sqrt{r_{n,d,q}}\), with constants that may depend on \(\alpha\), uniformly over the sample sizes, dimensions, and arrival floors satisfying its conditions. Alongside the squared-error order in \cref{obj:thm:matched-minimax-frontier}, this gives an uncertainty-quantification interpretation of the same comparison rate: its square root is the attainable and necessary order of maximal expected width for uniformly honest connected intervals. In particular, along sequences satisfying \cref{obj:ass:rare-arrival-slice}, the optimal expected length tends to zero exactly when the rate in \cref{obj:def:frontier-rate} tends to zero.
